% \titlepanel: the right of the title slide, which the closing slide repeats: a blue panel with a
% request typed at a prompt and, above it, the editor where the agent changes the code: one line
% removed (red), two added (green), the second still being written at a sparkle; a red edge on its
% left. Inside a full-page tikzpicture[x=1mm,y=1mm] whose origin is the page's top-left corner,
% y downwards negative; the panel bleeds 1 mm past the edges.
\newcommand\titlepanel{%
  \coordinate (pnw) at (102,0);
  \fill[udsblue] (102,1) rectangle (161,-91);
  % centred in the blue right of the red stripe
  \begin{scope}[shift={(pnw)},xshift=0.8mm]
    \fill[udsdark] (6,-15) rectangle (52,-60);
    \fill[udsdark!70!black] (6,-15) rectangle (52,-19.5);
    \foreach \i in {0,1,2} \fill[white!55!udsdark] (9+2.6*\i,-17.25) circle (0.65);
    % the agent's edit: rows 3 (removed) and 4 and 5 (added), marked in the gutter
    \fill[udsred!55!udsdark] (6,-31.4) rectangle (52,-34.6);
    \fill[testgreen!60!udsdark] (6,-34.6) rectangle (52,-41);
    \foreach \y/\m in {-33/$-$,-36.2/+,-39.4/+}
      \node[text=white,font=\footnotesize\bfseries,inner sep=0pt] at (9.5,\y) {\m};
    % code lines: row, indent level, segment widths; the first segment of a line is brighter
    \foreach \row/\ind/\segs in {0/0/{8,15},1/1/{5,3,12},2/1/{4,14},3/1/{6,17},4/1/{6,9,8},
        5/2/{8,7},6/2/{5,10},7/1/{11,6},8/1/{7,8},10/0/{9,5,11}} {
      \pgfmathsetmacro\y{-22.8-3.2*\row}
      \pgfmathsetmacro\x{13+3.5*\ind}
      \xdef\xs{\x}
      \foreach \w [count=\k] in \segs {
        \pgfmathsetmacro\mix{\k==1 ? 85 : 50}
        \fill[white!\mix!udsdark] (\xs,\y) rectangle ++(\w,-1.2);
        \pgfmathsetmacro\nx{\xs+\w+1.4}
        \xdef\xs{\nx}
      }
    }
    % the sparkle pair right of the line being written (row 5 ends at x = 36.4, row 4 at 42.3)
    \foreach \sx/\sy/\sr in {45.6/-39.2/2.5,47.6/-37.3/1.1}
      \fill[white,shift={(\sx,\sy)},scale=\sr] (0,1) .. controls (0,0.28) and (0.28,0) .. (1,0)
        .. controls (0.28,0) and (0,-0.28) .. (0,-1) .. controls (0,-0.28) and (-0.28,0) .. (-1,0)
        .. controls (-0.28,0) and (0,0.28) .. cycle;
    % the prompt: a chevron and three words, the first word in line with the code
    \draw[white,line width=0.9pt,rounded corners=3.5mm] (6,-64.5) rectangle (52,-71.5);
    \draw[udsred!45!white,line width=1.2pt] (8.9,-66.6) -- (10.5,-68) -- (8.9,-69.4);
    \foreach \x/\w in {13/9,23.4/4,28.8/12} \fill[white!85!udsblue] (\x,-67.4) rectangle ++(\w,-1.2);
  \end{scope}
  \fill[udsred] (102,1) rectangle (103.6,-91);}

% ------------------------------------------------------------------ title
\begin{frame}[plain,noframenumbering]
% Drawn inside the frame rather than over the page, so the speaker notes' thumbnail shows it too.
% The picture's origin is the page's top-left corner: a plain frame's body starts 8 mm in and
% 2 mm down.
\vbox to 0pt{\vskip-2mm\hbox to 0pt{\hskip-8mm\begin{tikzpicture}[x=1mm,y=1mm]
  \useasboundingbox (0,0) rectangle (160,-90);
  \titlepanel
  \node[anchor=north west,inner sep=0pt] at (8,-8) {\includegraphics[width=30mm]{\udslogo}};
  \node[anchor=north west,inner sep=0pt,text width=88mm,align=left] at (8,-36) {%
    {\fontsize{25}{29}\selectfont\bfseries\color{udsblue}\inserttitle\par}\vskip2mm
    {\fontsize{15}{18}\selectfont\color{udsdark}\insertsubtitle\par}\vskip4mm
    \textcolor{udsred}{\rule{14mm}{0.8pt}}\par\vskip3mm
    {\fontsize{11}{13}\selectfont\color{udsdark}\insertauthor\par}};
  \node[anchor=south west,inner sep=0pt] (d3m) at (8,-83) {\includegraphics[width=10mm]{assets/d3m.png}};
  \node[anchor=west,inner sep=0pt,text width=74mm,align=left,font=\fontsize{7}{8.8}\selectfont,
    execute at begin node={\hyphenpenalty=10000}] at ([xshift=3mm]d3m.east) {%
    {\color{udsblue}\insertinstitute\par}
    {\color{udsgrey}\insertdate\par}};
\end{tikzpicture}\hss}\vss}
% a frame with no tall content drops out of the notes' thumbnail; this invisible strut keeps it
\par\rule{0pt}{80mm}
\note{\textbf{0:00--0:20}\par
    Practical talk: how we prepare a repository for coding agents, how we hand over a task, and
      how we check the result.}
\end{frame}

% ---------------------------------------------------------------- outline
\begin{frame}{Outline}
  % each section's title as on its divider, and the divider's question, shortened to one line
  \small\renewcommand\arraystretch{1.32}\setlength\tabcolsep{4pt}
  \begin{tabularx}{\textwidth}{@{}>{\leavevmode\color{udsblue}\bfseries}p{4mm}>{\bfseries}p{61mm}X@{}}
    \toprule
    1 & From asking questions to delegating work & What must we provide, and what must we verify?\\
    2 & Environments and code checks             & Where does our code run, and what checks it?\\
    3 & Preparing the repository                 & What can the repository tell the agent?\\
    4 & Architecture Decision Records (ADRs)     & Why does the project work this way?\\
    5 & Version control with git                 & How do we keep and review every change?\\
    6 & Tests first                              & How do tests keep an agent on track?\\
    7 & Skills and an end-to-end workflow        & How does a request become a reviewed change?\\
    8 & A realistic starting point               & What should I change in my repository first?\\
    \bottomrule
  \end{tabularx}
  \note{\textbf{0:20--1:00}\par
    One running task carries the talk: splitting a segmentation dataset by specimen.\par
    The eight sections answer eight practical questions.\par
    Sections 2 and 5 are the groundwork that an agent needs in any project: an environment with
      automatic checks, and version control.\par
    A thread runs from the first section to the fourth: the better the repository is prepared,
      the shorter the request can be.\par
    The repository and its code are illustrative; the micrographs come from our HCCI SEM
      dataset.\par
    Claude Code, Codex and Matt Pocock's skills are the documented examples; the research-code
      examples are ours.\par
    Tool behaviour changes quickly; the appendix lists versions and access dates.}
\end{frame}
